% Source: https://artofproblemsolving.com/community/c3h1583331_geometry_countdown_help
% Stable identity: aops-topic-1583331-q13
Adjacent vertices are $2$ apart, so neighbouring circles are tangent and the pieces inside the decagon do not overlap. (Non-adjacent vertices are farther apart, since the decagon's diagonals are longer than its sides.) Near each vertex, the part of the circle inside the decagon is a sector whose angle is the interior angle of the decagon,
\[
\dfrac{(10-2)\cdot180^\circ}{10}=144^\circ .
\]

\begin{asy}
usepackage("amsmath");
size(6cm);
defaultpen(fontsize(10pt));
int n=10;real R=1/sin(pi/10);
pair[] V;for(int i=0;i<n;++i) V.push(R*dir(90+36*i));
for(int i=0;i<n;++i){
  pair p=V[i],q=V[(i+1)%n],r=V[(i+n-1)%n];
  real a1=degrees(q-p),a2=degrees(r-p);
  if(a2<a1) a2+=360;
  fill(p--arc(p,1,a1,a2)--cycle,paleblue);
}
for(int i=0;i<n;++i) draw(circle(V[i],1),gray);
draw(V[0]--V[1]--V[2]--V[3]--V[4]--V[5]--V[6]--V[7]--V[8]--V[9]--cycle,linewidth(1));
for(int i=0;i<n;++i) dot(V[i]);
label("$2$",(V[5]+V[6])/2+0.6*dir(288));
draw(arc(V[5],0.45,degrees(V[6]-V[5]),degrees(V[4]-V[5])+360),red);
label("$144^\circ$",V[5]+(0,0.72),red+fontsize(8pt));
label("\mbox{}",(min(currentpicture,user=true)+max(currentpicture,user=true))/2,black); // keeps the text after the figure black when LaTeX typesets the labels inline
\end{asy}

The ten sectors together have total angle $10\cdot144^\circ=1440^\circ=4\cdot360^\circ$, i.e.\ they make up four whole circles of radius $1$. The total area is
\[
4\cdot\pi\cdot1^2=\boxed{4\pi}\text{ square units}.
\]
